\documentclass[a4paper,12pt]{article}

\usepackage[T1]{fontenc}
\usepackage[spanish,provide=*]{babel}
\usepackage{geometry}
\usepackage{setspace}
\usepackage{hyperref}
\usepackage{listings}

\geometry{left=3cm,right=2.5cm,top=3cm,bottom=2.5cm}
\onehalfspacing
\hypersetup{colorlinks=true,urlcolor=blue,linkcolor=black,citecolor=black}

\lstset{
  language=C++,
  basicstyle=\ttfamily\small,
  frame=single,
  breaklines=true,
  numbers=left
}

\begin{document}

\title{Sistema Hospitalario Inteligente en C++}
\author{Proyecto ADA}
\date{2025}
\maketitle

\tableofcontents
\newpage

\section{Resumen}

El proyecto implementa un sistema inteligente de atencion hospitalaria en C++,
aplicando estructuras de datos y algoritmos estudiados en ADA: colas de prioridad,
tablas hash, ordenamiento, busqueda binaria, algoritmos voraces y benchmarking.

\section{Estructura del Codigo}

\begin{itemize}
  \item \texttt{main.cpp}: menu principal.
  \item \texttt{include/Paciente.h}: estructura de paciente y comparador.
  \item \texttt{include/Hospital.h}: clase principal del sistema.
  \item \texttt{src/Algoritmos.cpp}: QuickSort, MergeSort y busqueda binaria.
  \item \texttt{src/HistorialMedico.cpp}: historial medico con hash table.
  \item \texttt{src/GestorMedicos.cpp}: asignacion greedy de medicos.
  \item \texttt{src/Benchmark.cpp}: medicion empirica de tiempos.
\end{itemize}

\section{Compilacion}

\begin{lstlisting}
g++ main.cpp src/*.cpp -Iinclude -o hospital -std=c++17
./hospital
\end{lstlisting}

\section{Conclusiones}

La seleccion adecuada de estructuras y algoritmos mejora la eficiencia del sistema:
la cola de prioridad gestiona emergencias en $O(\log n)$, el historial usa busqueda
promedio $O(1)$ y los ordenamientos eficientes trabajan en $O(n \log n)$.

\end{document}
